\documentclass[12pt]{article}
\usepackage[utf8]{inputenc}
\usepackage[T1]{fontenc}
\usepackage[margin=1in]{geometry}
\usepackage{setspace}
\usepackage{xcolor}
\usepackage[colorlinks=true, urlcolor=blue!60!black, linkcolor=blue!60!black]{hyperref}
\setlength{\parindent}{0em}
\setlength{\parskip}{0.6em}
\renewcommand{\baselinestretch}{1.0}
\begin{document}
\begin{center}
{\LARGE \textit{Lemon v.\ Kurtzman}, 403 U.S.\ 602 (1971)\par}
\vspace{0.5em}
{\large Case Brief --- Two-Page Edition\par}
\end{center}

\section{Memory Jogger}
State salary and service subsidies for parochial schools fostered excessive entanglement with religion; the case announced the three-prong \textit{Lemon} test, abandoned in 2022 (\href{https://supreme.justia.com/cases/federal/us/403/602/}{Justia}).

\section{Detailed Case Facts}
Rhode Island paid salary supplements to teachers of secular subjects in nonpublic schools, nearly all Catholic, who promised not to teach religion. Pennsylvania reimbursed nonpublic schools for teachers' salaries, textbooks, and materials in four secular subjects, subject to state audits.

\section{Procedural History}
A Rhode Island district court struck down its Act for ``excessive entanglement''; a Pennsylvania district court dismissed the challenge. The Supreme Court consolidated the appeals.

\section{Judicial Votes}
Pennsylvania: 8--0, reversed. Rhode Island: 8--1, affirmed. Burger wrote for the Court; Douglas (with Black) and Brennan concurred; White concurred in part and dissented in part.

\section{Holding}
Both statutes violate the Establishment Clause. A law must have a secular purpose, a primary effect that neither advances nor inhibits religion, and no ``excessive government entanglement with religion.'' The Court decided only the entanglement prong.

\section{Analysis of Opinions}
\textbf{Majority:} Unlike a textbook, a teacher in a pervasively religious school cannot be inspected once; enforcing the secular-only rules would require continuing surveillance and audits, and annual appropriations would breed ``political division along religious lines.'' \textbf{Douglas and Brennan:} subsidies free church funds for religious teaching and use religious means for secular ends. \textbf{White:} the Court created an ``insoluble paradox,'' striking down both unsafeguarded aid and the safeguards.

\section{Examples: Future Applications}
\textbf{Same-Side:} state salary subsidies for teachers in religious schools, policed by inspectors. \textbf{Opposite-Side:} neutral vouchers directed by private choice (\textit{Zelman}). \textbf{Fence-Sitter:} public remedial tutors on parochial campuses, later allowed in \textit{Agostini}.

\section{Critique}
Justice Scalia compared \textit{Lemon} to a horror-movie ghoul that keeps rising from its grave; originalists called it ahistorical. Separationists defend its protection of religious minorities.

\section{Key Quotations}
\begin{quote}\sloppy
``First, the statute must have a secular legislative purpose; second, its principal or primary effect must be one that neither advances nor inhibits religion, \textit{Board of Education v.\ Allen}, 392 U.S.\ 236, 243 (1968); finally, the statute must not foster `an excessive government entanglement with religion.'\,''
\end{quote}
\noindent --- Chief Justice Burger, for the Court.

\end{document}
